% !TEX root = ../main.tex

% ==========================================
% KAPITEL 4: UMSETZUNG
% ==========================================
\chapter{Umsetzung}
Das vorherige Kapitel definiert das theoretische Konzept. Dieses Kapitel beschreibt die praktische Realisierung des Prototyps.

\section{Implementierung der Komponenten}

\subsection{Konstruktion und 3D-Druck}
\label{sec:3ddruck_umsetzung}
% --- PLATZHALTER MANUEL (3D-DRUCK UMSETZUNG) ---
[Hier fügt Manuel seinen sachlichen Text zur CAD-Modellierung in Fusion 360, den Maßen, dem Klappmechanismus und dem Druckvorgang ein.]

\subsection{Schaltungsdesign und Energiemanagement}
\label{sec:elektronik_umsetzung}
% --- PLATZHALTER FABIAN (ELEKTRONIK UMSETZUNG) ---
[Hier fügt Fabian seine Schaltpläne, die genaue Verkabelung (TP4056, Nicla) und die Entladungskurven der Akkus ein.]

\subsection{Entwicklung des Machine-Learning-Modells (Objekterkennung)}
\label{sec:ml_training}
% --- PLATZHALTER FABIAN (EDGE IMPULSE / ROBOFLOW) ---
[Hier fügt Fabian seinen Text zum Training des Modells mit Roboflow, Kaggle und Edge Impulse ein.]

\subsection{Software-Architektur und Datenverarbeitung}
\label{sec:software_implementierung}
Die softwareseitige Umsetzung erfordert die Koordination eines ressourcenbeschränkten Embedded-Systems (MicroPython) und eines mobilen Betriebssystems (Android). Die folgenden Abschnitte erläutern die Lösungsansätze für die Netzwerkkommunikation, die asynchrone Datenverarbeitung und das Speichermanagement detailliert.

% --- BLOCK 1: NETZWERK & APP-INFRASTRUKTUR ---

\subsubsection{Android Lifecycle-Management und Systementkopplung}
Da das Smartphone während der Nutzung der Assistenzbrille typischerweise in der Tasche des Nutzers verbleibt und der Bildschirm deaktiviert ist, neigt das Android-Betriebssystem dazu, Hintergrundprozesse aus Energiespargründen zu terminieren (Doze-Mode). Die Architektur lagert daher die gesamte Netzwerk- und Bluetooth-Logik in einen Android \texttt{Service} (\texttt{AscentaService.kt}) aus. 

Der Aufruf von \texttt{startForeground()} signalisiert dem System höchste Priorität. Ein weiteres kritisches Problem stellt der Standby-Modus dar, der den Empfang von UDP-Broadcast-Paketen auf Hardware-Ebene blockiert. Die Applikation löst dies durch die Akquirierung eines \texttt{WifiManager.MulticastLock}. Dieser Lock zwingt den WiFi-Chip, Broadcast-Pakete auch bei deaktiviertem Display durchzureichen:

\begin{lstlisting}[language=Java, caption={Akquirierung des MulticastLocks und Start des Foreground Services}]
val wifiManager = applicationContext.getSystemService(Context.WIFI_SERVICE) as WifiManager
multicastLock = wifiManager.createMulticastLock("AscentaUDP")
multicastLock?.setReferenceCounted(true)
multicastLock?.acquire()

startForeground(1, createNotification("Ascenta Active"))
\end{lstlisting}

\subsubsection{Asynchrone BLE-GATT State-Machine}
Das System überträgt die WLAN-Credentials (Provisioning) initial über das Bluetooth Low Energy (BLE) Generic Attribute Profile (GATT) vom Smartphone auf die Brille. Da Bluetooth-Operationen in Android stark asynchron und Event-getrieben ablaufen, implementiert der Code eine komplexe State-Machine innerhalb eines \texttt{BluetoothGattCallback}.

Sobald der \texttt{AscentaService} die Brille über den BLE-Scanner findet, baut er eine GATT-Verbindung auf. Das Betriebssystem meldet Statusänderungen an die Methode \texttt{onConnectionStateChange}. Bei erfolgreicher Verbindung initiiert die Applikation die \textit{Service Discovery}. Erst wenn die Services des Nicla Vision im Android-System registriert sind, schreibt die Applikation die SSID als Bytearray in die spezifische Characteristic. Ein erfolgreicher Schreibvorgang triggert \texttt{onCharacteristicWrite}. Daraufhin überträgt das System kaskadierend das Passwort und terminiert die BLE-Verbindung sauber:

\begin{lstlisting}[language=Java, caption={Ausschnitt der asynchronen GATT Callback-Struktur in Kotlin}]
override fun onServicesDiscovered(g: BluetoothGatt, s: Int) {
    val svc = g.getService(PROV_SERVICE_UUID)
    val ssidChar = svc?.getCharacteristic(SSID_CHAR_UUID)
    if (ssidChar != null) {
        ssidChar.value = ssid.toByteArray()
        g.writeCharacteristic(ssidChar) // Triggert asynchronen Schreibvorgang
    }
}

override fun onCharacteristicWrite(g: BluetoothGatt, char: BluetoothGattCharacteristic, status: Int) {
    if (status == BluetoothGatt.GATT_SUCCESS && char.uuid == SSID_CHAR_UUID) {
        val passChar = g.getService(PROV_SERVICE_UUID).getCharacteristic(PASS_CHAR_UUID)
        Thread.sleep(100) // Hardware-Pufferung des Nicla Vision abwarten
        passChar.value = pendingPass!!.toByteArray()
        g.writeCharacteristic(passChar)
    }
}
\end{lstlisting}

\subsubsection{Hybrides Netzwerkprotokoll (UDP \& TCP)}
Nach erfolgreichem BLE-Provisioning verbindet sich der Mikrocontroller mit dem WLAN und sendet periodisch \texttt{NICLA\_READY}-Broadcasts via UDP auf Port 5006. Eine Kotlin-Coroutine in der App lauscht asynchron auf diese Pakete und ermittelt die dynamisch vergebene IP-Adresse der Brille. 

\begin{lstlisting}[language=Java, caption={Asynchroner UDP-Discovery-Listener in Kotlin}]
udpSocket = DatagramSocket(UDP_DISC_PORT).apply { broadcast = true }
val packet = DatagramPacket(buffer, buffer.size)

while (isActive && !isConnected) {
    udpSocket?.receive(packet)
    val str = String(packet.data, 0, packet.length).trim()
    if (str == "NICLA_READY") {
        udpSocket?.close()
        connectTcp(packet.address) // TCP-Handshake einleiten
        break
    }
}
\end{lstlisting}

Nach der IP-Ermittlung etabliert die App den Hauptdatenstrom über eine verbindungsorientierte TCP-Verbindung (Port 5005). Die Architektur nutzt hierbei TCP anstelle von UDP, da dieses Protokoll Paketverluste bei der Übertragung großer binärer Payloads (WAV-Audio-Chunks oder JPEG-Bilder) strikt verhindert. Die Firmware konfiguriert den Socket auf dem Nicla Vision mittels \texttt{setblocking(False)}. Dies verhindert, dass abgebrochene Verbindungen den Mikrocontroller in einen Deadlock versetzen.

\subsubsection{Stream-Multiplexing und Custom Protocol Parser}
Die Nutzung einer einzigen TCP-Verbindung (Port 5005) für völlig unterschiedliche Datentypen kennzeichnet eine architektonische Besonderheit des Systems. Der Datenstrom (\textit{Stream}) interleavt kontinuierliche Textnachrichten (JSON-Telemetrie, Heartbeat-Pongs) mit großen binären Payloads (JPEG-Bilder, WAV-Audiodaten). 

Ein eigenes, leichtgewichtiges Applikationsprotokoll (Application Layer) verhindert eine Vermischung der Datentypen. Der Code stellt binären Übertragungen einen definierten Präfix inklusive Payload-Größe voran (z.\,B. \texttt{IMG\_START:12500}). Der Android \texttt{AscentaService} implementiert einen maßgeschneiderten Parser. Dieser Parser liest den eingehenden \texttt{InputStream} zunächst zeichenweise (Byte für Byte) bis zum Line-Feed (\texttt{\textbackslash n}) aus. Erkennt der Parser einen Binär-Präfix, schaltet er den Lesemodus dynamisch um, extrahiert exakt die angegebene Anzahl an Bytes in ein \texttt{ByteArray} und wechselt danach wieder in den Text-Modus:

\begin{lstlisting}[language=Java, caption={Dynamisches Stream-Multiplexing zur Trennung von Text- und Binärdaten}]
// Auslesen des TCP-InputStreams in der Coroutine
val line = readLineFromStream(input) // Zeichenweises Lesen bis \n

if (line.startsWith("IMG_START:")) {
    val size = line.substringAfter(":").toInt()
    val bytes = readBytesFromStream(input, size) // Blockweises Lesen
    val bmp = BitmapFactory.decodeByteArray(bytes, 0, bytes.size)
    if (bmp != null) handler.post { callback?.onImageReceived(bmp) }
    
} else if (line.startsWith("AUD_START:")) {
    val size = line.substringAfter(":").toInt()
    val bytes = readBytesFromStream(input, size)
    saveAudioForVosk(bytes)
    
} else if (line.startsWith("{")) {
    // Standard JSON-Parsing fuer Sensordaten
}
\end{lstlisting}

\subsubsection{Ausfallsicherheit durch Heartbeat-Mechanismus}
Die \texttt{MainActivity} implementiert einen asynchronen Watchdog, um Verbindungsabbrüche zuverlässig zu erkennen. Die Applikation sendet alle zehn Sekunden einen Ping an die Brille. Bleibt der Response für mehr als 25 Sekunden aus, schließt die Software die Verbindung sicher und reaktiviert den UDP-Discovery-Modus (\textit{Self-Healing}):

\begin{lstlisting}[language=Java, caption={Heartbeat-Coroutine zur Überwachung der Verbindungsstabilität}]
heartbeatJob = CoroutineScope(Dispatchers.IO).launch {
    while (isActive) {
        if (isConnected) {
            sendCommand("ping")
            delay(10000)
            if (System.currentTimeMillis() - lastPingResponse > 25000) {
                withContext(Dispatchers.Main) {
                    isConnected = false
                    stopHeartbeat()
                    showHome() // Zurueck in den Discovery-Status
                }
            }
        } else { delay(5000) }
    }
}
\end{lstlisting}

\subsubsection{Concurrency und Thread-Sicherheit in Kotlin}
Die Android-Applikation führt netzwerkbasierte I/O-Operationen (TCP/UDP-Sockets) und rechenintensive Machine-Learning-Prozesse (Vosk) simultan aus. Dafür implementiert der Code ein striktes Thread-Dispatching mittels \textit{Coroutines}. Sämtliche Sockets und blockierende Lesevorgänge operieren ausschließlich im \texttt{Dispatchers.IO} Thread-Pool. Sobald die Coroutine ein Datenpaket vollständig empfängt, führt sie einen synchronisierten Kontextwechsel zurück auf den Main-Thread durch (\texttt{handler.post}). Der Main-Thread aktualisiert anschließend die Benutzeroberfläche oder steuert die TTS-Engine an:

\begin{lstlisting}[language=Java, caption={Sicherer Kontextwechsel zwischen IO-Thread und Main-Thread in Kotlin}]
scope.launch(Dispatchers.IO) { // Start im Hintergrund-Thread
    val input = tcpSocket!!.getInputStream()
    while (isActive && isConnected) {
        val line = readLineFromStream(input) // Blockierender I/O-Aufruf
        
        if (line.startsWith("{")) {
            val json = JSONObject(line)
            val label = json.optString("label")
            
            // Sicherer Wechsel auf den Main-Thread fuer das UI/Audio-Feedback
            handler.post { 
                callback?.onDetectionResult(label, conf, dist) 
            }
        }
    }
}
\end{lstlisting}

% --- BLOCK 2: KI & SENSORIK AUF DER BRILLE ---

\subsubsection{Ressourcenmanagement und Modell-Integration (Nicla Vision)}
Die Firmware des Nicla Vision (\texttt{Firmware.py}) operiert unter starken Speicherrestriktionen. Das System reserviert initial einen Notfall-Puffer (\texttt{micropython.alloc\_emergency\_exception\_buf(100)}), um Out-of-Memory (OOM) Exceptions während Laufzeit-Interrupts zu verhindern. Das Skript integriert das trainierte FOMO-Modell ressourceneffizient direkt im Framebuffer des Mikrocontrollers (\texttt{load\_to\_fb=True}). 

\subsubsection{Die Kernschleife der Umgebungserfassung}
Die Hauptmethode \texttt{run\_active\_detection} auf dem Mikrocontroller steuert die kontinuierliche Überwachung der Umgebung und die Peripherie. Diese Funktion minimiert Latenzen und vermeidet blockierende Aufrufe. 

In jedem Iterationsschritt erfasst der Sensor ein neues Kamerabild (\texttt{sensor.snapshot()}) und übergibt dieses an das geladene Machine-Learning-Modell (\texttt{net.predict()}). Parallel wertet das System die Distanz über den I2C-Bus des VL53L1X aus. Erkennt der Algorithmus ein valides Objekt mit einer Konfidenz über dem Schwellenwert, konstruiert die Firmware dynamisch einen JSON-String. Dieser String bündelt das Label, den Konfidenzwert sowie die gemessene Distanz. Anschließend streamt das System dieses Paket asynchron über den TCP-Socket:

\begin{lstlisting}[language=Python, caption={Extraktion der Kameradaten und asynchroner JSON-Telemetrie-Stream}]
img = sensor.snapshot()
distance_mm = 0
if self.tof:
    try: distance_mm = self.tof.read()
    except: pass

raw_output = self.net.predict([img])
# ... (Parsing der Heatmap, Bestimmung von best_label und best_conf) ...

if best_label:
    json_str = '{"label": "%s", "conf": "%.2f", "dist": "%d", "x": "%d", "y": "%d"}\n' % (
        best_label, best_conf, distance_mm, best_x, best_y
    )
    self.send_tcp_packet(json_str.encode())
    
gc.collect() # Proaktive Defragmentierung des Speichers
\end{lstlisting}

\subsubsection{Auswertung des 4D-Tensors und Sensorfusion}
Die Auswertung der Vorhersagedaten des FOMO-Modells stellt eine besondere algorithmische Herausforderung dar. Das quantisierte Modell generiert einen vierdimensionalen Tensor (Heatmap), der das Bild in ein Grid unterteilt. Die Firmware iteriert über das X/Y-Grid. Findet der Algorithmus ein lokales Maximum über dem Schwellenwert ($0,45$), berechnet die Firmware den absoluten Bildschwerpunkt (Zentroid) und fusioniert die 2D-Bildkoordinaten mit der Z-Achse (Distanz in Millimetern) des ToF-Sensors.

\subsubsection{Low-Level Hardware-Kommunikation (I2C \& Registerauslesung)}
Die Firmware liest den integrierten Fuel-Gauge-IC (MAX17262) über den I2C-Bus (Inter-Integrated Circuit) aus, um dem Nutzer zuverlässige Batteriewarnungen zu geben. Der Mikrocontroller adressiert den Chip unter der Hexadezimal-Adresse \texttt{0x36} und fordert zwei Bytes ab dem Register \texttt{0x06} (RepSOC) an. Da der Chip die Daten im Little-Endian-Format überträgt, konvertiert die Software die rohen Bytes in einen Integer und dividiert den Wert anschließend durch $256,0$. Dieses Verfahren berechnet den exakten prozentualen Ladezustand:

\begin{lstlisting}[language=Python, caption={Register-Level I2C-Auslesung des MAX17262 Fuel-Gauge-ICs}]
def get_soc(self):
    if not self.fuel_gauge_i2c: return None
    try:
        # Auslesen von 2 Bytes aus Register 0x06 (RepSOC)
        data = self.fuel_gauge_i2c.readfrom_mem(0x36, 0x06, 2)
        
        # Konvertierung von Little-Endian Bytes zu Integer
        soc_raw = int.from_bytes(data, 'little')
        
        # Umrechnung in Prozent gemaess Datenblatt
        soc = soc_raw / 256.0 
        return soc
    except:
        return None
\end{lstlisting}

% --- BLOCK 3: AUDIO & INTERAKTION ---

\subsubsection{Kinetische Interaktionssteuerung (Tap-Detection)}
Eine knopflose Bedienung aktiviert die Spracherkennung. Das System liest hierfür den Beschleunigungssensor (LSM6DSOX) via SPI aus. Anstatt absolute Beschleunigungswerte zu nutzen, berechnet die Firmware den Ruck ($\Delta g$) zwischen zwei Frame-Messungen. Dieser Filter ignoriert normale Kopfbewegungen:

\begin{lstlisting}[language=Python, caption={Berechnung von Delta-G zur Erkennung eines Double-Taps}]
x, y, z = self.lsm.accel()
dx = abs(x - self.last_accel[0])
dy = abs(y - self.last_accel[1])
dz = abs(z - self.last_accel[2])
self.last_accel = (x, y, z)

delta_g = dx + dy + dz

if delta_g > 1.75:
    now = time.ticks_ms()
    diff = time.ticks_diff(now, self.last_tap_time)
    
    if 100 < diff < 600:
        self.send_tcp_packet(b'{"action": "tap"}\n')
        self.req_audio = True 
        self.last_tap_time = 0
    elif diff > 600:
        self.last_tap_time = now
\end{lstlisting}

\subsubsection{Speichereffizientes Audio-Streaming}
Eine Audioaufnahme (2,5\,s, 16\,kHz, Mono) erzeugt Datenmengen, die den freien Arbeitsspeicher des Nicla Vision übersteigen, falls das System sie am Stück puffert. Die Firmware implementiert daher einen \textit{Circular Buffer} (\texttt{g\_audio\_chunks}). 

Ein Hardware-Callback (\texttt{\_audio\_callback}) schreibt die Audiodaten während der Aufnahme iterativ in vorab allozierte 1024-Byte-Chunks. Die Nutzung eines \texttt{memoryview} verhindert, dass der Interpreter bei der Array-Zuweisung neuen Speicher alloziert (Zero-Copy-Prinzip). Nach der Aufnahme generiert die Firmware dynamisch einen 44 Byte großen RIFF/WAVE-Header mittels \texttt{struct.pack} und streamt diesen direkt über den TCP-Socket:

\begin{lstlisting}[language=Python, caption={Speichereffizientes Befüllen des Circular Buffers mittels memoryview in MicroPython}]
def _audio_callback(self, buf):
    global g_audio_written
    if not self.recording: return
    
    mv = memoryview(buf)
    l = len(mv)
    off = 0
    
    while l > 0:
        if self.rec_idx >= len(g_audio_chunks): 
            self.recording = False
            break
            
        chunk = g_audio_chunks[self.rec_idx]
        space = AUDIO_CHUNK_SIZE - self.rec_offset
        amt = min(l, space)
        
        # Speicherzuweisung ohne RAM-Neuallokation (Zero-Copy)
        chunk[self.rec_offset : self.rec_offset+amt] = mv[off : off+amt]
        
        self.rec_offset += amt
        off += amt
        l -= amt
        g_audio_written += amt
        
        # Puffer-Index inkrementieren, wenn Chunk voll ist
        if self.rec_offset >= AUDIO_CHUNK_SIZE:
            self.rec_idx += 1
            self.rec_offset = 0
\end{lstlisting}

\subsubsection{Dynamisches Asset-Management des Akustikmodells}
Das File-Management stellt ein spezifisches technisches Hindernis bei der Nutzung der lokalen Vosk-Spracherkennung unter Android dar. Die zugrundeliegende Vosk-Engine erfordert absolute Dateipfade auf einem unkomprimierten Dateisystem. Die Applikation löst dieses Problem durch einen asynchronen Entpackungsprozess beim Erststart der App. In einer IO-Coroutine iteriert die Funktion \texttt{copyAssets} rekursiv durch das komprimierte Modell-Verzeichnis der APK und kopiert die binären Dateien in den internen \texttt{filesDir}-Speicherbereich des Smartphones:

\begin{lstlisting}[language=Java, caption={Rekursives, asynchrones Entpacken des Vosk-Modells in Kotlin}]
private fun copyAssets(path: String, out: String) {
    val list = assets.list(path) ?: return
    if (list.isNotEmpty()) {
        File(out).mkdirs()
        list.forEach { copyAssets("$path/$it", "$out/$it") }
    } else {
        val dir = File(out).parentFile
        if (dir != null && !dir.exists()) dir.mkdirs()
        assets.open(path).use { input -> 
            FileOutputStream(out).use { input.copyTo(it) } 
        }
    }
}
\end{lstlisting}

\subsubsection{Audio-Verarbeitung und kognitive Entlastung}
Auf dem Smartphone übergibt das System die resultierende WAV-Datei an die \textit{Vosk Offline Engine}. Die Applikation liest die Datei schrittweise in 4096-Byte-Blöcken ein, um UI-Freezes zu verhindern. Nach der Auswertung generiert das System Warnungen über die \texttt{TextToSpeech}-API. Die App implementiert eine dynamische Cooldown-Logik, um den blinden Nutzer vor einer auditorischen Überlastung zu bewahren: Objekte in unkritischer Entfernung (über 50\,cm) lösen eine Blockade von sechs Sekunden aus. Nähert sich ein Objekt jedoch in den kritischen Nahbereich (unter 50\,cm), reduziert der Algorithmus den Cooldown drastisch und generiert eine direkte \glqq Kollisionswarnung!\grqq.

% --- BLOCK 4: UI & DEBUGGING ---

\subsubsection{I/O-Auslagerung beim Bild-Streaming}
Die Brille bietet zur Verifikation des Sichtfeldes die Funktion, ein Live-Bild an die Smartphone-App zu senden. Das direkte Versenden eines unkomprimierten Bitmaps über TCP führt jedoch zu einer sofortigen OOM-Exception. 

Als Lösungsansatz nutzt die Firmware das Dateisystem als \textit{Swap-Speicher}. Das Skript schreibt das Kamerabild mit stark reduzierter Qualität (40\,\%) als JPEG in den lokalen Flash-Speicher (\texttt{temp.jpg}). Danach zerstört die Firmware die Bildinstanz im RAM manuell (\texttt{del img}) und ruft den Garbage Collector auf. Erst im Anschluss liest das System die komprimierte JPEG-Datei iterativ in kleinen 1024-Byte-Blöcken aus und streamt diese über das Netzwerk:

\begin{lstlisting}[language=Python, caption={Verwendung des Flash-Speichers als Swap zur Vermeidung von OOM-Exceptions}]
def process_image(self):
    path = "temp.jpg"
    try:
        img = sensor.snapshot()
        img.save(path, quality=40) # Speichern im Flash
        
        # Manuelle Zerstörung der Bildinstanz zur sofortigen RAM-Freigabe
        del img
        gc.collect()
        
        size = os.stat(path)[6]
        self.send_tcp_packet(f"IMG_START:{size}\n".encode())
        
        # Iteratives Auslesen und Senden der komprimierten Datei
        with open(path, 'rb') as f:
            while True:
                chunk = f.read(1024)
                if not chunk: break
                self.send_tcp_packet(chunk)
                
        uos.remove(path) # Temporäre Datei sicher löschen
    except: 
        pass
    finally:
        gc.collect()
\end{lstlisting}

\subsubsection{Benutzeroberfläche und App-Architektur (GUI)}
Die Architektur der Applikation folgt dem \textit{Single-Activity-Pattern}. Der Code unterteilt die eigentlichen Bildschirminhalte in modulare \texttt{Fragments}. Das \texttt{InfoFragment} visualisiert den aktuellen Verbindungsstatus sowie den via I2C ausgelesenen Akkustand der Hardware. Das \texttt{ImageFragment} empfängt Bilddatenpakete, dekodiert diese und rendert sie als Android \texttt{Bitmap}. Zudem implementiert dieses Fragment eine Funktion zur dauerhaften Speicherung von Umgebungsaufnahmen über die \texttt{MediaStore}-API des Betriebssystems (Scoped Storage). Das \texttt{AudioFragment} bietet eine Live-Transkription der Vosk-Engine. Dies ermöglicht Entwicklern, die Genauigkeit des akustischen Modells bei verschiedenen Umgebungsgeräuschen zu evaluieren.


\section{Herausforderungen in der Entwicklung}
Die Realisierung eines kompakten, tragbaren Assistenzsystems, welches rechenintensive Machine-Learning-Prozesse mit kontinuierlicher Netzwerkkommunikation vereint, bringt diverse interdisziplinäre Herausforderungen mit sich. Diese erfordern im Projektverlauf iterative Anpassungen an der Hardware-Architektur sowie tiefgreifende Software-Optimierungen.

\subsection{Mechanische und geometrische Restriktionen}
% --- PLATZHALTER MANUEL ---
[Hier beschreibt Manuel seine Probleme beim 3D-Druck, Scharniere etc. sachlich.]

\subsection{Energiemanagement im begrenzten Bauraum}
% --- PLATZHALTER FABIAN ---
[Hier beschreibt Fabian den Kompromiss mit den kleineren Akkus sachlich.]

\subsection{Ressourcenlimits und Speicher-Fragmentierung (Software)}
Auf softwareseitiger Ebene stellt der extrem begrenzte Arbeitsspeicher (RAM) des Nicla Vision die größte Hürde dar. Frühe Prototypen der Firmware stürzen regelmäßig mit \textit{Out-Of-Memory} (OOM) Exceptions ab. Dies hat zwei Hauptursachen:

Erstens übersteigt die Pufferung von Audiodaten den verfügbaren Speicher. Das System verzichtet daher auf Standard-Audio-Bibliotheken zugunsten eines selbst entwickelten \textit{Circular Buffers}. Dieser Buffer fragmentiert die Audiodaten über \texttt{memoryview}-Zuweisungen speichereffizient in 1024-Byte-Blöcke.

Zweitens verursacht die \textit{Heap-Fragmentierung} Abstürze während der Bildverarbeitung. MicroPython alloziert Speicher dynamisch. Das ständige Erzeugen und Verwerfen von Objekten (wie \texttt{image.snapshot()} bei 240x240 Pixeln) hinterlässt mikroskopische Lücken im Speicher. Benötigt das FOMO-Modell anschließend einen zusammenhängenden Speicherblock für den Tensor, schlägt die Allokation trotz ausreichendem Gesamtspeicher fehl. Die Firmware implementiert daher eine strikte, deterministische Garbage-Collection-Strategie. Der Code erzwingt den Befehl \texttt{gc.collect()} nach jedem Frame-Snapshot und in den \texttt{finally}-Blöcken der Netzwerk-Routinen, um den RAM proaktiv zu defragmentieren.

Auf Seiten der Android-Applikation resultieren die größten Herausforderungen aus den restriktiven Energiesparmaßnahmen moderner Android-Versionen (Doze-Mode). Da die Applikation primär im Hintergrund bei gesperrtem Smartphone-Display operiert, blockiert das Betriebssystem standardmäßig den Empfang von UDP-Broadcasts. Dies verhindert den \glqq Zero-Config\grqq-Verbindungsaufbau. Die Implementierung eines \texttt{Foreground Services} in Kombination mit einem expliziten \texttt{WifiManager.MulticastLock} zwingt die Hardware jedoch dazu, eintreffende Netzwerkpakete zuverlässig an die Applikation durchzureichen. Dieser Mechanismus stellt die Systemstabilität letztlich sicher.


\section{Tests und Validierung}
Eine Reihe von technischen Validierungen belegt die Praxistauglichkeit des \glqq Ascenta\grqq-Prototyps. Die Tests fokussieren sich auf die Kernanforderungen eines Assistenzsystems: Latenzzeit, Zuverlässigkeit der Sensorik und energetische Ausdauer.

\subsection{Netzwerklatenz und Streaming-Performance}
Die Messung der Systemlatenz ermittelt die Zeitdifferenz zwischen der physischen Platzierung eines Objekts im Sichtfeld der Kamera und der Generierung des ersten Audio-Warnsignals (TTS) auf dem Smartphone. Dies stellt ein kritisches Sicherheitsmerkmal für blinde Nutzer dar.

Unter optimalen Laborbedingungen (kurze Distanz zwischen Brille und Smartphone, wenig Interferenz im 2,4\,GHz WiFi-Band) ergibt sich eine durchschnittliche Latenzzeit von etwa [XX] Millisekunden für die Bilderkennung und den reinen Telemetrie-Datentransfer (JSON). Der Transfer eines $2,5$-sekündigen Audio-Streams (16\,kHz, Mono) via TCP an die Vosk-Engine inklusive anschließender Offline-Transkription dauert im Durchschnitt [XX] Sekunden. Diese Ergebnisse belegen, dass die gewählte \glqq Split-Brain\grqq-Architektur den Anforderungen an ein Echtzeitsystem gerecht wird.

\subsection{Auswertung der kinetischen Interaktion (Tap-Detection)}
Eine Versuchsreihe mit [XX] bewussten Eingaben (Taps) und [XX] simulierten Störbewegungen (schnelles Gehen, Kopfschütteln) evaluiert die Zuverlässigkeit der Double-Tap-Erkennung über die IMU (LSM6DSOX).

\begin{itemize}
    \item \textbf{True-Positive-Rate:} Bei [XX] von [XX] beabsichtigten Doppeltipps aktiviert sich die Sprachaufnahme korrekt (Erkennungsrate von [XX]\,\%).
    \item \textbf{False-Positive-Rate:} Während des simulierten Gehens über eine Distanz von [XX] Metern treten [XX] Fehlerkennungen auf. 
\end{itemize}
Der gewählte Schwellenwert für die Ableitung der Beschleunigung ($\Delta g > 1,75$) filtert normale Kopfbewegungen somit adäquat heraus und setzt dennoch keine unverhältnismäßig hohe Schlagkraft des Nutzers voraus.

\subsection{Energetische Evaluierung (Akkulaufzeit)}
Die kontinuierliche Auslesung des integrierten MAX17262 Fuel-Gauge-ICs ermittelt die Laufzeit des Gesamtsystems. Das Testverfahren prüft zwei Betriebsmodi bei der verbauten Gesamtkapazität von 1240\,mAh (Parallelbetrieb von zwei 620\,mAh LiPo-Zellen):

\textbf{Standby-Modus (Kamera deaktiviert, nur UDP-Heartbeat aktiv):} \\
In diesem Modus beschränkt sich die Leistungsaufnahme auf die Aufrechterhaltung der WiFi-Verbindung. Der gemessene Entladestrom liefert eine theoretische Maximallaufzeit von [XX] Stunden.

\textbf{Aktiver Modus (Kontinuierliche FOMO-Erkennung und TCP-Stream):} \\
Im aktiven Modus wertet das Machine-Learning-Modell kontinuierlich Snapshots aus und sendet Datenpakete an die Applikation. Unter dieser maximalen Auslastung sinkt die Laufzeit auf etwa [XX] Stunden und [XX] Minuten. Dieser praktikable Wert ermöglicht den Alltagseinsatz während Einkäufen oder Wegstrecken problemlos. Die Tests zeigen zudem keine signifikante thermische Drosselung (Thermal Throttling) des Cortex-M7 Prozessors.